\documentclass[12pt]{article}
\usepackage[utf8]{inputenc}
\usepackage[T1]{fontenc}
\usepackage{lmodern}
\usepackage{microtype}
\usepackage[hidelinks]{hyperref}
\usepackage[a4paper, margin=3cm]{geometry}
\usepackage{parskip}

\title{Sopori Concert - Welcome Address}
\date{22/05/26}


\begin{document}
\maketitle

Goede avond, en hartelijk welkom.
On behalf of Stichting Sahaja Yoga Nederland, it is my privilege to welcome you to today's special event.
My name is Ganesh and I will be guiding you through this evening.

Dear guests, today we have with us an exceptional musician, who is not only himself a pillar of
Indian Classical Music and the Santoor today, but also stems from a musical family that has shaped
the landscape since more than 10 generations. He and his father — famously called the Saint of
Santoor — are special guests for us not only because of their musical prowess, but also because of
their spiritual understanding. Their music cannot be considered as entertainment. It has to be
considered an offering — an offering to the divine and everything profound in this world. We are
therefore extremely delighted to have with us today Pt. Abhay Sopori, son of the legendary Pt.
Bhajan Sopori.

This evening promises to be an experience for the audience. An experience of stillness and peace. It
is something that increasingly becomes a scarcity in these times that we live in. It is also important
to me to acknowledge all current conflicts that are going on at the moment. I believe it is our
societal duty to not ignore those, despite the difficulty of keeping up with the vast and numbing
amount of news that reach us daily.

I mention them because I believe they immediately concern us and that we can play a part in
solving them. The message I want to give today is that we are all part of those conflicts — whether
we want to be or not. As psychologist Carl Jung has explained, there is a collective consciousness
that we all belong to. A consciousness that is made up of each one's individual consciousness. In
other words, it matters to the wider world whether we are peaceful within or not. Even though we
all condemn the ongoing global conflicts, we can, at the same time, think of numerous times where
we haven't been peaceful to the standards we set for ourselves. The grudges we held, the times we
hurt someone, the people we couldn't forgive. I think they all matter.

My point here is not to judge anyone or make us feel bad, but rather to encourage each one of us. If
our actions matter, wherever we may be, we are not helpless. On the contrary, we can directly shape
the world. That means if we are able to solve our inner conflicts — if we establish a peace within
ourselves — we can automatically propagate peace into the world.

And this is not a new idea. This knowledge exists since thousands of years — written down in
ancient scriptures. Be it the Vedas in the Indian tradition, the Quran in the Arabic world, the Bible in
the Middle East, the Tao in the Chinese tradition. They all speak of an energy that can be felt, and a
resurrection that takes place within.

That is the vision Dr. Nirmala Salve, also known as Shri Mataji Nirmala Devi, has for mankind. In
her lifetime she studied the behaviour of human beings and explored methods to solve their issues at
the root. Her efforts resulted in the method of Sahaja Yoga that facilitates human transformation
through self-realisation. That is the union with the inner self, also called Atma in Sanskrit, or Spirit
in English. The novelty of Shri Mataji's contribution lies in making this experience tangible and
accessible to the masses.

When the inner energy, called Kundalini, is awakened, she rises and enables the state of thoughtless
awareness. This is a state of inner silence that allows us to connect to the deeper self, the Spirit.
This state of awareness, and the connection to that self which is nothing but pure joy, results in a
state of absolute inner peace. Where reactions of the mind don't affect us. Where emotions don't
dominate us. Where we can truly be our self.

And just like technological advancements are evident throughout evolution, the spiritual evolution
is taking place in the same way. In my opinion these are times to celebrate, and it gives me
tremendous joy to share that experience with all of you.

We will begin with an Aarti in honour of Shri Mataji — that is also an invitation for the divine to be
present with us in this room tonight. This will be followed by a brief guided meditation. And then
we will welcome the artists onto the stage, and allow them to take us on a journey into the depths of
our being.
I invite you all to rise.

\end{document}
